% Generated draft: requires mathematical and source review.
\documentclass[11pt]{article}
\usepackage[margin=25mm]{geometry}
\usepackage{fontspec}
\setmainfont{Latin Modern Roman}
\usepackage{amsmath,amssymb,mathtools,unicode-math}
\setmathfont{Latin Modern Math}
\usepackage{enumitem,needspace,xurl,hyperref}
\setlength{\parindent}{0pt}
\setlength{\parskip}{5pt}
\setlength{\emergencystretch}{3em}
\newcommand{\N}{\mathbb{N}}
\newcommand{\Z}{\mathbb{Z}}
\newcommand{\Q}{\mathbb{Q}}
\newcommand{\R}{\mathbb{R}}
\newcommand{\C}{\mathbb{C}}
\newcommand{\F}{\mathbb{F}}
\newcommand{\A}{\mathbb{A}}
\newcommand{\PP}{\mathbb{P}}
\newcommand{\E}{\mathbb{E}}
\newcommand{\eps}{\varepsilon}
\newcommand{\dd}{\,\mathrm{d}}
\begin{document}
\section*{UnsolvedMath ID: 4600015}
\textbf{Source problem code:} catalogue entry
\subsection{Definitions and mathematical statement}
Let $A$ be a finite nonempty alphabet, $X=A^{\Z_{\geq0}}$ with the product topology, and let $\sigma_+:X\to X$ be the left shift. An automorphism of this one-sided full shift is a homeomorphism $T:X\to X$ satisfying $T\sigma_+=\sigma_+T$. Although $\sigma_+$ is generally not invertible, the automorphism $T$ is invertible.

For a compatible metric $d$ on $X$, $T$ is expansive if
\[
\exists\delta>0\ \forall x,y\in X,\qquad
\left[d(T^n x,T^n y)<\delta\text{ for all }n\in\Z\right]\Longrightarrow x=y.
\]
The source asks whether, for every such expansive $T$, there are a finite nonempty alphabet $B$ and a homeomorphism $h:X\to B^{\Z}$ such that
\[
h\circ T=\sigma_B\circ h,
\]
where $\sigma_B$ is the two-sided full shift. The size of $B$ is not specified. This is conjugacy of the dynamical system $(X,T)$, not of the noninvertible map $\sigma_+$.

The supplied research summary attributes to Nasu the weaker conclusion that an expansive automorphism of a one-sided shift of finite type is conjugate to a two-sided shift of finite type. A shift of finite type is defined by finitely many forbidden finite blocks; it need not be a full shift. Therefore that quoted theorem does not by itself establish the source's full-shift conclusion. The summary alternates between full shift and shift of finite type, so its solved label cannot settle that distinction without a precise theorem and any additional hypotheses.
\subsection{Short English statement}
Is every expansive invertible map commuting with a one-sided full shift conjugate to a two-sided full shift? The supplied solved-status summary quotes a shift-of-finite-type conclusion, which is weaker than the full-shift target.
\subsection{Sources}
\begin{itemize}
\item Local UnsolvedMath registry, permanent ID 4600015, exact source code catalogue entry, ``One-sided full-shift automorphisms''; proposed year 2008, no proposer specified. Registry status solved, classification SOLVED-IN-LITERATURE.
\item Mike Boyle, \emph{Open Problems in Symbolic Dynamics} (2008), Problem/Question/Conjecture 12.2, \url{https://www.math.umd.edu/~mboyle/papers/openfinalsub3nov2007.pdf}. Supplied pdf-text extraction accessed July 29, 2026. Historical-source provenance says it is presented as open and no later authoritative resolution check is recorded. No PDF was accessed here.
\item The supplied research summary attributes to Nasu the statement ``Every expansive automorphism of a one-sided SFT is conjugate to a two-sided SFT'' and says a Boyle talk lists 12.2 as solved. It supplies no precise Nasu publication or talk reference, and also uses stronger full-shift language. These claims are recorded with their differing scope rather than conflated.
\item Provenance flags status and rights NEEDS\_REVIEW: public access was verified, but redistribution permission was not established. The claimed later solution is a supplied assessment requiring source and theorem-scope verification.
\end{itemize}
\end{document}
